\section{Additional details on Bayesian network learning (Learn)}
\label{sec:S5}

\subsection{Construction of the learning dataset}
\label{sec:S5.1}

Each version of a crisis, that is, a crisis played under a given management policy, provides one observation; the main dataset contains 8,000 of them. The variables are organized according to the causes that a diagnosis must be able to distinguish: the exposure of the network, the disruption, the management policy and the actual response of the operations manager, supplemented by the mechanism of the drop and by the resilience obtained (Table~5 of the main manuscript).

The reference demand cannot be used as a context variable, because it is calibrated on the volume that can be delivered during the incident and therefore depends on the disruption itself. It is replaced by the target position in the topology. A site or a link is a bottleneck if its removal cuts the customer off from all sources. On the ISOMORPH network, the Nashville hub, the Atlanta warehouse and the link between them are bottlenecks; congestion and demand surges affect the whole network. This variable depends only on the structure of the network and carries over to any network under study. The policy is represented by the profile alone, because its components (threshold, delay, pace, order of levers) are fully redundant with it in a dataset where they are fixed by profile. An absent response, when the service level never fell below the critical threshold, is coded by an explicit ``not applicable'' state.

The parameters of the curve are kept as intermediate variables, because they describe the mechanism through which a cause produces its effect (Table~\ref{tab:S3}). Three are discarded for lack of variation in the dataset: $k$, close to 1 by construction; $q$, close to 0 because the crises resolve before the end of the observed period; and $g$, almost always zero. Among the indicators, IPC, which faithfully reflects the lost service and has an unambiguous direction, is the prognosis target; TRP and CRD complement it. Three indicators are discarded: SRA, redundant with the parameters; ERR, too weakly dispersed; and PR, whose favorable direction depends on the depth of the drop.

All versions are kept, including those with an imperfect fit, because discarding poorly fitted crises would mainly eliminate short crises, and thus part of the successes of the reactive profile. Continuous variables are discretized into three equal-frequency classes, whose thresholds are computed on the training sample and then frozen for new crises.

\subsection{Constrained learning}
\label{sec:S5.2}

Learning is performed with the BayesArtist tool \citep{bayesartist}. The structure is learned by greedy hill-climbing local search with the Bayesian information criterion (BIC) \citep{schwarz1978}. The order of the tiers is imposed as a list of forbidden arcs, with at most four parents per variable, so that no learned structure can violate the temporal order of the crisis; this combination of domain knowledge and data follows the principle of \citet{heckerman1995}. The probability tables are estimated by counting with Laplace smoothing, which avoids assigning zero probability to a combination of causes absent from the dataset. Inference is exact, by variable elimination or junction tree \citep{koller2009}.

The policy, an intervention compared on strictly identical crises, has an effect on resilience that is identifiable without adjustment. The levers actually mobilized, in contrast, depend on what the operations manager perceived and act as mediators.

Since BayesArtist is not publicly distributed, the same learning can be carried out with any Bayesian network software that learns structure and parameters from a CSV file and accepts forbidden arcs, for example Hugin, GeNIe, the R package bnlearn or the Python library pgmpy; as search algorithms differ from one tool to another, slight differences in structure are possible.

\subsection{Inference, recommendation and validation protocol}
\label{sec:S5.3}

In prognosis, the network gives the distribution of resilience for a context and a policy, for example the probability of a low cumulative loss when a severe storage site closure is managed by a prudent profile. In diagnosis, it goes back from a degraded resilience to the most probable causes: a bottleneck hit, a disruption that is hard to overcome, late perception, decision inertia, or unsuitable levers.

Recommendation combines the two. Given a new crisis (fully or partly described) and a resilience objective, each candidate policy is evaluated by its probability of reaching this objective, which can be read directly from the network since the policy is an intervention. The rule adopted favors sobriety: among the policies whose probability exceeds a threshold $\theta$ set by the decision maker, the least interventionist one is recommended, the policies being ordered by the average number of levers they engage (no reaction, delayed, prudent, reactive). If none reaches the threshold, the most probable policy is recommended. This rule reflects the fact that every lever has a cost that the model does not value; for equal resilience, mobilizing fewer levers is preferable.

The same question can be asked backward: set the objective (for example a low cumulative loss) and the known context, then read the distribution of the profile. In general, going back from an effect to a decision only yields an association, because the observed decisions depend on the situation. Here, the profile is a root variable, assigned according to a uniform distribution independently of the context. Its posterior probability is therefore proportional to its probability of reaching the objective. Backward inference ranks the policies exactly as prognosis does. With one query per context, it provides a recommendation map, which identifies the most effective policy; the sobriety rule remains necessary to select the most economical one among those that are sufficient. This property does not hold for the levers, which depend on what the operations manager perceives.

Validation exploits a rare property of the dataset, since each crisis is observed under all four policies. The 2,000 scenarios are randomly split into 1,600 training scenarios and 400 test scenarios; the split is made on scenarios, as the four versions of the same crisis are not independent. Prognosis is assessed on the prediction of the cumulative loss class from only the variables known at decision time, using accuracy and the Brier score, compared with the majority class and with a model based on the policy alone. Diagnosis is assessed by comparing, for high-loss crises, the posterior distribution of the causes with their observed frequency. Recommendation is assessed directly. Since each test crisis was played under all policies, we check whether the recommended policy actually reached the objective and measure the share of recommendations that are more sober than the reactive profile. A learning curve, from 200 to 1,600 scenarios, shows whether the size of the dataset is sufficient.
